Learning in artificial life must change a continuing body's responses through experience. Useful acquired content, however, need not outperform a simple response. We study these distinctions through paid local consequences, interventions on acquired content and strong physiological controls in explicitly separated cellular models. In our primary conservative-body assay, policies receive only locally completed scalar energy transfers and physically transported internal records, without observed transport displacements, reward gradients or field coordinates. A standard perturbation-response rule acquires useful action in both tested conditions, passes reward reversal in one, and produces a later unreinforced aftereffect in both. Clearing the acquired response leaves one aftereffect successful, whereas rotating it leaves neither successful. When only execution changes, continued estimation and an initialized fixed run-and-tumble response both succeed; the fixed response spends less on motion. Central and material-carried history controls support retained choices, including unseen gaps, without evidence that the tested recurrent-gain plasticity adds value. An equivalent fixed-material readout also reproduces learned mechanical responses. In an externally supplied birth/death material model, learned conditional maintenance reduces injury under equal protection purchases, yet fails to dominate immediate physiology after a relationship change. Stable-relation timing interventions reveal small benefits of prepared reserves. Together, these results establish bounded uses of acquired content and identify the supplied information and physical alternatives that explain their limits. Learning, adaptive response and material renewal remain distinct capabilities rather than an already unified regenerating learner.
